% !TEX program = xelatex
% Статья по материалам НИР/ВКР «Выявление кода, сгенерированного большими языковыми моделями»
% Сборка: tectonic article.tex  или  xelatex article.tex && xelatex article.tex

\documentclass{misisarticle}

\usepackage{tikz}
\usetikzlibrary{positioning, arrows.meta}

% --- Классификационные коды ---
\grnti{28.23.37, 28.23.15}
\asjc{1702, 1712}
\oecd{1.02}

% --- Название ---
\title{Многоканальные карты вероятностей токенов для выявления программного кода, сгенерированного большими языковыми моделями}

% --- Автор ---
\misisauthor{А.~И.~Беляев}{кафедра инженерной кибернетики НИТУ «МИСиС»}{Москва, Россия}{m2302576@edu.misis.ru}

% --- Аннотация и ключевые слова ---
\abstract{%
  Рассматривается задача бинарной классификации фрагментов исходного кода: требуется определить, написан ли фрагмент человеком или сгенерирован большой языковой моделью. За основу взят метод CodeVision, который раскладывает лог-вероятности токенов в двумерную карту, повторяющую расположение символов в тексте программы, и классифицирует её свёрточной сетью. Метод использует один канал и зависит от проприетарного программного интерфейса, поэтому плохо воспроизводится и уязвим к простым преобразованиям кода. Предлагается модификация с тремя изменениями. Карта становится многоканальной: к лог-вероятности добавляются энтропия распределения следующего токена и ранг наблюдаемого токена, которые вычисляет открытая языковая модель. Оценка свёрточной сети объединяется со статистическими детекторами (средняя лог-вероятность, отношение лог-вероятности к лог-рангу, оценка Binoculars) через логистическую мета-модель, обученную на внефолдовых прогнозах. Сеть обучается с регуляризацией согласованности на семантически эквивалентных преобразованиях кода: переформатировании, переименовании идентификаторов, изменении комментариев. Приводится математическая постановка задачи и описание методов. Протокол проверки включает метрики AUC и TPR при FPR не выше 1~\%, бутстреп-интервалы, а также перенос на языки программирования и модели-генераторы, не участвовавшие в обучении. Данные состоят из решений Project CodeNet и кода трёх открытых генераторов.
}

\keywords{%
  CodeVision, энтропия и ранг токена, Binoculars, стекинг детекторов, регуляризация согласованности, Project CodeNet
}

\begin{document}

\maketitle

% ============================================================
\section{Введение}

Языковые модели для программирования (Codex, StarCoder, DeepSeek-Coder, Qwen2.5-Coder) пишут синтаксически корректный и стилистически правдоподобный код по короткому текстовому заданию~\cite{qwen25coder, deepseekcoder}. Классический антиплагиат сравнивает фрагменты между собой и ищет источник заимствования. Против кода, сгенерированного с нуля, он бессилен: источника нет. Проблема авторства затрагивает образование (самостоятельность выполнения заданий), разработку (лицензионная чистота кодовой базы) и информационную безопасность.

Автоматическое выявление машинного кода опирается на то, что модель-генератор оставляет в тексте статистический след: выбираемые ею токены в среднем более вероятны и менее «неожиданны», чем токены, выбранные человеком. Метод CodeVision~\cite{codevision} переводит этот след в изображение: лог-вероятности токенов раскладываются в двумерную карту по позициям символов, и карта классифицируется моделью компьютерного зрения. Двумерная раскладка сохраняет отступы, переносы строк и блочную структуру программы, которые теряются в одномерной последовательности.

У исходного метода три ограничения. Вероятности берутся из проприетарного программного интерфейса, поэтому результаты дорого и трудно воспроизвести. Карта содержит один канал, хотя распределение следующего токена несёт больше информации, чем вероятность одного выбранного токена. Устойчивость к преобразованиям, которые не меняют смысла программы (переименование, переформатирование), почти не изучена, хотя именно такие правки делает пользователь, желающий скрыть происхождение кода~\cite{codegptsensor}.

Цель работы~--- построить модификацию CodeVision на открытой языковой модели, которая точнее исходного метода, устойчива к семантически эквивалентным преобразованиям кода и переносится на модели-генераторы, не встречавшиеся при обучении.

Вклад работы:
\begin{itemize}
  \item многоканальная карта, в которой к лог-вероятности добавлены энтропия распределения следующего токена и ранг наблюдаемого токена;
  \item объединение визуального классификатора со статистическими детекторами, включая оценку Binoculars~\cite{binoculars}, через мета-модель, обученную на внефолдовых прогнозах;
  \item обучение с регуляризацией согласованности на семантически эквивалентных преобразованиях кода;
  \item протокол проверки с доверительными интервалами и долей обнаружения при низкой доле ложных срабатываний, включающий перенос на новые языки, новые генераторы и смену опорной модели.
\end{itemize}

% ============================================================
\section{Анализ работ}

Методы выявления текста и кода, созданного языковыми моделями, делятся на пять групп.

\textit{Статистические методы без обучения} вычисляют по опорной модели характеристики текста и сравнивают их с порогом. Простейшие признаки~--- средняя лог-вероятность, ранг и энтропия токенов. DetectGPT~\cite{detectgpt} сравнивает правдоподобие текста и его случайных возмущений: машинный текст лежит вблизи локального максимума правдоподобия, и возмущения его заметно понижают. Fast-DetectGPT~\cite{fastdetectgpt} заменяет возмущения аналитической условной кривизной, что ускоряет метод на два порядка. DetectLLM~\cite{detectllm} использует отношение лог-вероятности к лог-рангу. Binoculars~\cite{binoculars} делит перплексию текста под одной моделью на перекрёстную перплексию между двумя родственными моделями и тем самым убирает зависимость от того, насколько текст «сложен» сам по себе. Для кода эти методы теряют точность: низкая энтропия программ оставляет мало места для различий~\cite{zeroshotcode}.

\textit{Обученные классификаторы} строятся поверх предобученных моделей кода. GPTSniffer~\cite{gptsniffer} дообучает CodeBERT и точен на своём распределении, но, по данным авторов, плохо работает, когда обучающие и тестовые данные собраны из независимых источников.

\textit{Цифровые водяные знаки} встраивают в генерацию скрытый сигнал, смещая выбор токенов~\cite{kirchenbauer}. Вариант SWEET~\cite{sweet} применяет смещение только к токенам с высокой энтропией, чтобы не ломать код. Водяной знак работает лишь для модели, владелец которой его внедрил, и не помогает против открытых моделей.

\textit{Методы возмущений и переписывания для кода} учитывают его синтаксис. DetectCodeGPT~\cite{detectcodegpt} возмущает код вставкой пробелов и переносов строк вместо переписывания внешней моделью. Другой подход просит модель переписать фрагмент и измеряет, насколько мало изменился результат: машинный код меняется меньше~\cite{rewriting}. Оба подхода требуют многократных проходов опорной модели на каждый фрагмент.

\textit{Перцептуальные методы} представлены CodeVision~\cite{codevision}, который сохраняет пространственную структуру кода и классифицирует карту лог-вероятностей моделями ViT~\cite{vit} и ResNet~\cite{resnet}. Авторы сообщают о высокой точности на данных Project CodeNet~\cite{codenet} для нескольких языков, но используют один канал и закрытый интерфейс. Устойчивость к изменениям кода в этой группе не изучалась. В работе CodeGPTSensor+~\cite{codegptsensor} показано, что обученные детекторы теряют точность на коде с переименованными идентификаторами и изменённой структурой, а состязательное обучение на таких примерах эту потерю устраняет.

Ни одна группа не решает задачу целиком. Статистические методы обобщаются, но чувствительны к порогу и длине фрагмента. Обученные классификаторы точны, но привязаны к распределению обучения. Перцептуальный подход сочетает сильные стороны обеих групп: признаки статистические, а решающее правило обучаемое. Поэтому он выбран в качестве основы, а его ограничения (один канал, закрытый интерфейс, неизученная устойчивость) определяют направления модификации.

% ============================================================
\section{Постановка задачи}

Дана выборка $D = \{(x_j, y_j)\}_{j=1}^{N}$, где $x_j$~--- фрагмент исходного кода, $y_j \in \{0, 1\}$~--- метка: $0$~--- код написан человеком, $1$~--- сгенерирован языковой моделью из множества генераторов $\mathcal{G}$. Каждый фрагмент относится к одному из языков программирования $\mathcal{L}$, машинный~--- ещё и к генератору $g_j \in \mathcal{G}$.

Требуется построить оценку $f\colon \mathcal{X} \to [0, 1]$ вероятности $P(y = 1 \mid x)$ и решающее правило
\begin{equation}
  \hat{y}(x) = \mathbb{I}\left[\, f(x) \ge \tau \,\right],
  \label{eq:rule}
\end{equation}
где $\mathbb{I}[\cdot]$~--- индикатор, порог $\tau$ выбирается на валидационной выборке, не пересекающейся с тестовой.

Оценка ищется в виде композиции
\begin{equation}
  f(x) = h\bigl(s_w(\Phi(x)),\, z(x)\bigr),
  \label{eq:composition}
\end{equation}
где $\Phi\colon \mathcal{X} \to \mathbb{R}^{C \times H \times W}$~--- отображение кода в $C$-канальную карту, $s_w$~--- свёрточная сеть с параметрами $w$, $z(x) \in \mathbb{R}^{K}$~--- вектор скалярных статистик, $h$~--- мета-модель. Отображение $\Phi$ и статистики $z$ вычисляются по фиксированной опорной языковой модели и не обучаются.

Основной критерий качества~--- площадь под ROC-кривой
\begin{equation}
  \mathrm{AUC}(f) = P\bigl(f(x^{+}) > f(x^{-})\bigr),
  \label{eq:auc}
\end{equation}
где $x^{+}$ и $x^{-}$~--- случайные машинный и человеческий фрагменты. Ложное обвинение человека дороже пропуска машинного кода, поэтому дополнительно используется доля обнаружения при ограниченной доле ложных срабатываний:
\begin{equation}
  \mathrm{TPR}_{\alpha}(f) = \max_{\tau}\ \mathrm{TPR}(\tau) \quad \text{при} \quad \mathrm{FPR}(\tau) \le \alpha,
  \label{eq:tpr}
\end{equation}
где $\alpha = 0{,}01$.

К детектору предъявляются три требования.
\begin{enumerate}
  \item \textit{Точность}: $\mathrm{AUC}$ и $\mathrm{TPR}_{\alpha}$ модификации не ниже, чем у одноканального CodeVision на тех же данных.
  \item \textit{Устойчивость}: для семантически эквивалентного преобразования $T$ падение $\Delta_T = \mathrm{AUC}(f; D_{\text{test}}) - \mathrm{AUC}(f; T(D_{\text{test}}))$ близко к нулю.
  \item \textit{Обобщение}: детектор, обученный на $D \setminus D_g$ без кода генератора $g$, сохраняет высокое $\mathrm{AUC}$ на коде $g$; то же для отложенного языка $l \in \mathcal{L}$ и для смены опорной модели при тестировании.
\end{enumerate}

% ============================================================
\section{Описание методов}

Схема метода показана на рис.~\ref{fig:pipeline}. Фрагмент кода токенизируется и проходит через опорную модель, по распределениям следующего токена строится многоканальная карта и скалярные статистики, затем свёрточная сеть и мета-модель выдают итоговую оценку.

\begin{figure}[htbp]
  \centering
  \begin{tikzpicture}[
      node distance=3.5mm,
      box/.style={draw, rounded corners=1pt, align=center, font=\footnotesize,
                  minimum width=0.56\columnwidth, inner sep=2.5pt},
      side/.style={draw, rounded corners=1pt, align=center, font=\footnotesize,
                   text width=0.25\columnwidth, inner sep=2.5pt},
      arr/.style={-{Stealth[length=1.6mm]}}]
    \node[box] (code) {Фрагмент кода $x$};
    \node[box, below=of code] (lm) {Опорная модель: $p_i(\cdot)$, $i = 1, \dots, n$};
    \node[box, below=of lm] (map) {Карта $\Phi(x)$: каналы $\ell$, $H$, $r$};
    \node[box, below=of map] (cnn) {Свёрточная сеть $s_w(\Phi(x))$};
    \node[box, below=of cnn] (meta) {Мета-модель $h(s_w, z)$, порог $\tau$};
    \node[side, right=5mm of map] (stats) {Статистики $z(x)$};
    \node[side, right=5mm of cnn] (aug) {Обучение на $T(x)$};
    \draw[arr] (code) -- (lm);
    \draw[arr] (lm) -- (map);
    \draw[arr] (map) -- (cnn);
    \draw[arr] (cnn) -- (meta);
    \draw[arr] (lm.east) -| (stats.north);
    \draw[arr] (stats.east) -- ++(2mm,0) |- (meta.east);
    \draw[arr, dashed] (aug.west) -- (cnn.east);
  \end{tikzpicture}
  \caption{Схема предлагаемого детектора}
  \label{fig:pipeline}
\end{figure}

\subsection{Признаки токенов}

Фрагмент $x$ токенизируется в последовательность $t_1, \dots, t_n$. Опорная модель с открытыми весами задаёт распределение следующего токена $p_i(v) = p(v \mid t_1, \dots, t_{i-1})$ по словарю $V$. В отличие от закрытого интерфейса, она даёт доступ ко всему распределению, что и делает возможными признаки, кроме вероятности одного токена. В работе используется модель Qwen2.5-Coder-0.5B-Instruct~\cite{qwen25coder}.

Для каждой позиции $i$ вычисляются лог-вероятность, ранг и энтропия:
\begin{align}
  \ell_i &= \log p_i(t_i), \label{eq:logprob}\\
  r_i &= \bigl|\{v \in V : p_i(v) > p_i(t_i)\}\bigr|, \label{eq:rank}\\
  H_i &= -\sum_{v \in V} p_i(v) \log p_i(v). \label{eq:entropy}
\end{align}
Лог-вероятность~\eqref{eq:logprob}~--- единственный канал исходного CodeVision. Ранг~\eqref{eq:rank} не зависит от того, насколько «острым» было распределение, и показывает, выбрала бы модель этот токен сама. Энтропия~\eqref{eq:entropy} описывает неопределённость модели в данной позиции, а не выбранный токен: место, где возможен один вариант, перестаёт быть свидетельством против человека.

Для сравнения реализованы и другие каналы (табл.~\ref{tab:channels}): разрыв двух наибольших вероятностей, лексический тип токена и условная кривизна Fast-DetectGPT~\cite{fastdetectgpt}
\begin{equation}
  c_i = \frac{\ell_i - \mu_i}{\sigma_i + \varepsilon}, \quad \mu_i = -H_i,
  \label{eq:curv}
\end{equation}
где $\sigma_i^2 = \sum_{v} p_i(v)\,(\log p_i(v))^2 - \mu_i^2$~--- дисперсия лог-вероятности токена, выбранного из $p_i$. Набор каналов определяется абляцией: каждый канал добавляется или исключается при прочих равных условиях.

\begin{table}[htbp]
  \caption{Каналы карты}
  \label{tab:channels}
  \centering
  \begin{tabular}{>{\raggedright\arraybackslash}p{0.3\columnwidth}>{\raggedright\arraybackslash}p{0.6\columnwidth}}
    \toprule
    \textbf{Канал} & \textbf{Смысл} \\
    \midrule
    $\ell_i$ & лог-вероятность токена (база CodeVision) \\
    $r_i$ & ранг токена в распределении \\
    $H_i$ & энтропия распределения \\
    $p_i^{(1)} - p_i^{(2)}$ & разрыв двух наибольших вероятностей \\
    тип токена & ключевое слово, идентификатор, литерал и~т.\,д. \\
    $c_i$ & условная кривизна~\eqref{eq:curv} \\
    \bottomrule
  \end{tabular}
\end{table}

\subsection{Двумерное представление}

Отображение $\Phi$ сохраняет раскладку исходного текста. Символ в строке $a$ и столбце $b$ занимает ячейку $(a, b)$ карты размера $H \times W$, и в неё записывается вектор признаков токена, покрывающего этот символ. Пустые ячейки (за концом строки и за последней строкой) получают нулевое значение, фрагменты длиннее $H$ строк или $W$ символов обрезаются. Каждый канал нормируется по формуле $(u - m_k)/s_k$, где $m_k$ и $s_k$~--- среднее и стандартное отклонение канала $k$ по занятым ячейкам обучающей выборки. Параметры нормировки фиксируются после обучения и применяются к валидационным и тестовым картам без пересчёта.

\subsection{Классификатор и регуляризация согласованности}

Карта подаётся в свёрточную сеть ResNet~\cite{resnet} с $C$ входными каналами и глобальным усредняющим пулингом, выход которой $s_w(\Phi(x)) \in [0, 1]$. Для устойчивости к правкам кода используется набор семантически эквивалентных преобразований $\mathcal{T}$: переформатирование (отступы, пробелы, переносы), согласованное переименование идентификаторов, вставка и удаление комментариев. Для обучающего примера $x_j$ случайно выбирается $T \in \mathcal{T}$ и строится карта $\tilde{M}_j = \Phi(T(x_j))$.

Параметры сети находятся минимизацией функционала
\begin{multline}
  L(w) = \frac{1}{N}\sum_{j=1}^{N} \beta_{y_j}\,\mathrm{CE}\bigl(s_w(M_j), y_j\bigr) \\
  + \lambda\, \frac{1}{N}\sum_{j=1}^{N} \mathrm{KL}_{\mathrm{sym}}\bigl(s_w(M_j) \,\|\, s_w(\tilde{M}_j)\bigr),
  \label{eq:loss}
\end{multline}
где $M_j = \Phi(x_j)$, $\mathrm{CE}$~--- перекрёстная энтропия, $\beta_c = N / (2 N_c)$~--- вес класса $c$ при $N_c$ примерах в классе, $\lambda$~--- вес регуляризатора, $\mathrm{KL}_{\mathrm{sym}}(p\,\|\,q) = \mathrm{KL}(p\,\|\,q) + \mathrm{KL}(q\,\|\,p)$ для бернуллиевских распределений. Первое слагаемое учит отличать классы, второе штрафует за изменение ответа при преобразовании кода. Это аналог состязательной регуляризации TRADES~\cite{trades} с той разницей, что возмущения берутся из множества правок кода, а не из шара в пространстве признаков. Преобразования применяются только при обучении. При оценке устойчивости они же служат атаками на тестовую выборку.

\subsection{Объединение со статистическими детекторами}

Свёрточная сеть видит локальные паттерны, а статистики без обучения~--- глобальные свойства фрагмента. Вектор $z(x)$ включает среднюю лог-вероятность $\bar{\ell} = \frac{1}{n}\sum_i \ell_i$, отношение лог-вероятности к лог-рангу $\mathrm{LRR} = \bar{\ell} \big/ \frac{1}{n}\sum_i \log(1 + r_i)$ по DetectLLM~\cite{detectllm} и другие простые статистики карты. Если доступна вторая модель с тем же токенизатором, добавляется оценка Binoculars~\cite{binoculars}
\begin{equation}
  B(x) = \frac{-\frac{1}{n}\sum_i \log p^{P}_i(t_i)}{\frac{1}{n}\sum_i \sum_{v} p^{P}_i(v)\bigl(-\log p^{O}_i(v)\bigr)},
  \label{eq:binoculars}
\end{equation}
где $p^{P}$ и $p^{O}$~--- распределения двух родственных моделей (в работе инструктивная и базовая версии Qwen2.5-Coder-0.5B). Знаменатель~--- перекрёстная энтропия моделей~--- показывает, насколько «удивителен» текст для модели в среднем. Деление на него убирает вклад сложности самого фрагмента.

Мета-модель~--- логистическая регрессия
\begin{equation}
  h(x) = \sigma\bigl(\gamma_0 + \gamma_1\, \hat{s}(x) + \boldsymbol{\gamma}^{\top} z(x)\bigr),
  \label{eq:meta}
\end{equation}
где $\sigma$~--- логистическая функция. Чтобы мета-модель не переобучилась на уверенных ответах сети по уже виденным примерам, используется стекинг~\cite{stacking}. Объединение обучающей и валидационной частей делится на $k$ блоков, сеть обучается на $k-1$ из них, а оценка $\hat{s}(x)$ для объектов оставшейся части берётся из модели, которая их не видела.

\subsection{Данные и протокол оценки}

Человеческий код~--- 220 решений задач из Project CodeNet~\cite{codenet}, по 44 на языки Python, C, C++, Java и Go. Машинный код~--- 200 фрагментов, полученных автором от трёх открытых моделей по банку из 31 алгоритмической задачи: Qwen2.5-Coder-0.5B-Instruct (150 фрагментов), Qwen2.5-Coder-1.5B-Instruct~\cite{qwen25coder} и DeepSeek-Coder-1.3B-Instruct~\cite{deepseekcoder} (по 25). Все генераторы дают равное число фрагментов на каждый язык, поэтому язык не выдаёт метку.

Используются три разбиения. Случайное делит выборку в пропорции 60/15/25~\% на обучающую, валидационную и тестовую части. В разбиении по языкам один язык целиком откладывается в тест. В разбиении по генераторам из обучения исключается весь код одной модели, и детектор проверяется на нём и на человеческом коде. Устойчивость к опорной модели проверяется заменой модели, которая строит тестовые карты.

Порог~\eqref{eq:rule} выбирается на валидации по максимуму сбалансированной точности $\frac{1}{2}(\mathrm{TPR} + \mathrm{TNR})$. Для $\mathrm{AUC}$~\eqref{eq:auc} и $\mathrm{TPR}_{\alpha}$~\eqref{eq:tpr} строятся 95-процентные бутстреп-интервалы~\cite{bootstrap}. Тестовая часть одного разбиения содержит около 50 человеческих фрагментов, а при $\alpha = 0{,}01$ одна ошибка уже превышает допустимую долю ложных срабатываний. Поэтому $\mathrm{TPR}_{\alpha}$ считается по объединённым оценкам нескольких запусков с разными начальными значениями генератора случайных чисел. Конфигурации сравниваются по перекрытию интервалов, а не по точечным значениям.

% ============================================================
\begin{thebibliography}{99}

\bibitem{qwen25coder} Hui, B. Qwen2.5-Coder Technical Report / B. Hui, J. Yang, Z. Cui [et al.]. -- arXiv:2409.12186, 2024.

\bibitem{deepseekcoder} Guo, D. DeepSeek-Coder: When the Large Language Model Meets Programming~-- The Rise of Code Intelligence / D. Guo, Q. Zhu, D. Yang [et al.]. -- arXiv:2401.14196, 2024.

\bibitem{codevision} Xu, Z. CodeVision: Detecting LLM-Generated Code Using 2D Token Probability Maps and Vision Models / Z. Xu, V. S. Sheng // Proceedings of the AAAI Conference on Artificial Intelligence. -- 2025. -- Vol.~39. -- arXiv:2501.03288.

\bibitem{codegptsensor} Yin, X. Detecting LLM-generated Code with Subtle Modification by Adversarial Training / X. Yin, X. Li, C. Ni, X. Xu, X. Yang. -- arXiv:2507.13123, 2025.

\bibitem{binoculars} Hans, A. Spotting LLMs With Binoculars: Zero-Shot Detection of Machine-Generated Text / A. Hans, A. Schwarzschild, V. Cherepanova [et al.] // Proceedings of the 41st International Conference on Machine Learning. -- PMLR, 2024. -- Vol.~235. -- P.~17519--17537.

\bibitem{detectgpt} Mitchell, E. DetectGPT: Zero-Shot Machine-Generated Text Detection using Probability Curvature / E. Mitchell, Y. Lee, A. Khazatsky, C. D. Manning, C. Finn // Proceedings of the 40th International Conference on Machine Learning. -- PMLR, 2023. -- Vol.~202. -- P.~24950--24962.

\bibitem{fastdetectgpt} Bao, G. Fast-DetectGPT: Efficient Zero-Shot Detection of Machine-Generated Text via Conditional Probability Curvature / G. Bao, Y. Zhao, Z. Teng, L. Yang, Y. Zhang // International Conference on Learning Representations (ICLR). -- 2024. -- arXiv:2310.05130.

\bibitem{detectllm} Su, J. DetectLLM: Leveraging Log Rank Information for Zero-Shot Detection of Machine-Generated Text / J. Su, T. Y. Zhuo, D. Wang, P. Nakov // Findings of the Association for Computational Linguistics: EMNLP 2023. -- 2023. -- P.~12395--12412.

\bibitem{zeroshotcode} Yang, X. Zero-Shot Detection of Machine-Generated Codes / X. Yang, K. Zhang, H. Chen [et al.]. -- arXiv:2310.05103, 2023.

\bibitem{gptsniffer} Nguyen, P. T. GPTSniffer: A CodeBERT-based classifier to detect source code written by ChatGPT / P. T. Nguyen, J. Di Rocco, C. Di Sipio [et al.] // Journal of Systems and Software. -- 2024. -- Vol.~214. -- Art.~112059. -- DOI 10.1016/j.jss.2024.112059.

\bibitem{kirchenbauer} Kirchenbauer, J. A Watermark for Large Language Models / J. Kirchenbauer, J. Geiping, Y. Wen [et al.] // Proceedings of the 40th International Conference on Machine Learning. -- PMLR, 2023. -- Vol.~202. -- P.~17061--17084.

\bibitem{sweet} Lee, T. Who Wrote this Code? Watermarking for Code Generation / T. Lee, S. Hong, J. Ahn [et al.]. -- arXiv:2305.15060, 2023.

\bibitem{detectcodegpt} Shi, Y. Between Lines of Code: Unraveling the Distinct Patterns of Machine and Human Programmers / Y. Shi, H. Zhang, C. Wan, X. Gu // IEEE/ACM 47th International Conference on Software Engineering (ICSE). -- 2025. -- arXiv:2401.06461.

\bibitem{rewriting} Ye, T. Uncovering LLM-Generated Code: A Zero-Shot Synthetic Code Detector via Code Rewriting / T. Ye, Y. Du, T. Ma [et al.] // Proceedings of the AAAI Conference on Artificial Intelligence. -- 2025. -- arXiv:2405.16133.

\bibitem{vit} Dosovitskiy, A. An Image is Worth 16x16 Words: Transformers for Image Recognition at Scale / A. Dosovitskiy, L. Beyer, A. Kolesnikov [et al.] // International Conference on Learning Representations (ICLR). -- 2021. -- arXiv:2010.11929.

\bibitem{resnet} He, K. Deep Residual Learning for Image Recognition / K. He, X. Zhang, S. Ren, J. Sun // Proceedings of the IEEE Conference on Computer Vision and Pattern Recognition (CVPR). -- 2016. -- P.~770--778.

\bibitem{codenet} Puri, R. Project CodeNet: A Large-Scale AI for Code Dataset for Learning a Diversity of Coding Tasks / R. Puri, D. S. Kung, G. Janssen [et al.] // NeurIPS Datasets and Benchmarks Track. -- 2021. -- arXiv:2105.12655.

\bibitem{trades} Zhang, H. Theoretically Principled Trade-off between Robustness and Accuracy / H. Zhang, Y. Yu, J. Jiao [et al.] // Proceedings of the 36th International Conference on Machine Learning. -- PMLR, 2019. -- Vol.~97. -- P.~7472--7482.

\bibitem{stacking} Wolpert, D. H. Stacked generalization / D. H. Wolpert // Neural Networks. -- 1992. -- Vol.~5, № 2. -- P.~241--259.

\bibitem{bootstrap} Efron, B. An Introduction to the Bootstrap / B. Efron, R. J. Tibshirani. -- New York : Chapman \& Hall, 1993. -- 436~p.

\end{thebibliography}

\end{document}
